% =================================================
% Теоретический лист. Урок 1
% =================================================

\begin{center}
  {\fontsize{25}{30}\selectfont\bfseries\sffamily\color{Primary}
    Цифры, числа и разряды}

  \vspace{0.3em}

  {\large\sffamily\color{GrayText}
    Урок 1 из 3. Обозначение натуральных чисел}
\end{center}

\vspace{0.45em}

\section{Числа и цифры}

При счёте предметов мы называем числа:
\[
  1,\ 2,\ 3,\ 4,\ldots
\]
Такие числа называют \textbf{натуральными}. Для записи чисел используют
всего десять знаков --- \textbf{цифр}.

\digitcardsfigure

\begin{propertybox}[Цифра и число --- не одно и то же]
  \textbf{Цифра} --- знак, используемый при записи числа.
  \textbf{Число} показывает количество или порядок предметов.

  Например, число \(507\) записано тремя цифрами: \(5\), \(0\) и \(7\).
\end{propertybox}

Одну и ту же цифру можно использовать в записи разных чисел:
\[
  5,\qquad 15,\qquad 50,\qquad 505.
\]
Число, записанное одной цифрой, называют \textbf{однозначным}; число,
записанное несколькими цифрами, --- \textbf{многозначным}.

\begin{notebox}
  Запись числа не начинают с нуля. Поэтому \(057\) --- не новая запись
  трёхзначного числа: это то же число \(57\).
\end{notebox}

\section{Значение цифры зависит от её места}

Сравним четыре записи:
\[
  6,\qquad 60,\qquad 600,\qquad 6\,000.
\]
Во всех них встречается цифра \(6\), но она обозначает соответственно
\(6\) единиц, \(6\) десятков, \(6\) сотен и \(6\) тысяч.

\begin{propertybox}[Позиционный принцип]
  Место цифры в записи числа называют \textbf{разрядом}. Значение цифры
  зависит от того, в каком разряде она стоит.
\end{propertybox}

Разряды считают справа налево:
\[
  \text{единицы},\quad
  \text{десятки},\quad
  \text{сотни},\quad
  \text{единицы тысяч},\quad
  \text{десятки тысяч},\ldots
\]

\placevaluefigure

Для числа \(58\,342\) получаем:

\begin{center}
  \renewcommand{\arraystretch}{1.35}
  \begin{tabularx}{0.92\textwidth}{@{}>{\bfseries\sffamily}X>{\centering\arraybackslash}p{0.2\textwidth}>{\centering\arraybackslash}p{0.25\textwidth}@{}}
    \toprule
    \rowcolor{TableHeader}
    Разряд & Цифра & Значение цифры \\
    \midrule
    десятки тысяч & \(5\) & \(50\,000\) \\
    единицы тысяч & \(8\) & \(8\,000\) \\
    сотни & \(3\) & \(300\) \\
    десятки & \(4\) & \(40\) \\
    единицы & \(2\) & \(2\) \\
    \bottomrule
  \end{tabularx}
\end{center}

\section{Зачем в записи числа нужны нули}

\zeroplacefigure

В числе \(40\,508\):
\[
  4\text{ десятка тысяч},\quad
  0\text{ тысяч},\quad
  5\text{ сотен},\quad
  0\text{ десятков},\quad
  8\text{ единиц}.
\]

\begin{warningbox}[Ноль нельзя просто удалить]
  Если убрать нули из записи \(40\,508\), получится число \(458\).
  Цифры \(4\), \(5\) и \(8\) переместятся в другие разряды, поэтому
  значение числа изменится.
\end{warningbox}

\section{Разложение числа по разрядам}

Каждое многозначное число можно представить как сумму значений его цифр.

\expandedformfigure

Полная запись выглядит так:
\[
  47\,305=
  4\cdot10\,000+7\cdot1\,000+3\cdot100+0\cdot10+5.
\]
Слагаемое \(0\cdot10\) обычно не пишут:
\[
  47\,305=40\,000+7\,000+300+5.
\]

\begin{trainingtask}[Два направления]
  \begin{enumerate}
    \item Разложите число:
      \[
        63\,204=60\,000+3\,000+200+4.
      \]
    \item Соберите число:
      \[
        70\,000+800+6=70\,806.
      \]
  \end{enumerate}
\end{trainingtask}

\begin{warningbox}[Типичная ошибка]
  Неверно:
  \[
    50\,204=50\,000+2\,000+4.
  \]
  Цифра \(2\) стоит в разряде сотен, поэтому верно:
  \[
    50\,204=50\,000+200+4.
  \]
\end{warningbox}

\section{Цифра разряда и общее число разрядных единиц}

В числе \(347\):
\begin{itemize}
  \item цифра десятков равна \(4\);
  \item всего полных десятков --- \(34\), потому что \(347=34\cdot10+7\).
\end{itemize}

\begin{notebox}[Важно различать формулировки]
  Вопросы «какая \textbf{цифра} стоит в разряде десятков?» и
  «сколько всего \textbf{десятков} содержится в числе?» имеют разные ответы.
\end{notebox}

\section{Короткий алгоритм}

\begin{conclusionbox}[Как разобрать многозначное число]
  \begin{enumerate}
    \item Начинайте с крайней правой цифры: это разряд единиц.
    \item Двигайтесь влево: десятки, сотни, тысячи и далее.
    \item Назовите цифру каждого разряда.
    \item Чтобы разложить число, замените каждую ненулевую цифру её значением.
  \end{enumerate}
\end{conclusionbox}

\begin{practicebox}[Проверьте себя]
  \begin{enumerate}
    \item Сколько цифр в записи числа \(70\,604\)?
    \item Какая цифра стоит в разряде сотен числа \(47\,382\)?
    \item Каково значение цифры \(6\) в числе \(26\,415\)?
    \item Соберите число: \(40\,000+2\,000+50+9\).
  \end{enumerate}
\end{practicebox}

\begin{notebox}[Ответы]
  1) \(5\); \quad 2) \(3\); \quad 3) \(6\,000\); \quad 4) \(42\,059\).
\end{notebox}

\begin{questionsbox}
  \begin{enumerate}
    \item Чем число отличается от цифры?
    \item Почему одна и та же цифра может иметь разное значение?
    \item Какую роль играет ноль внутри записи числа?
    \item Как разложить число на разрядные слагаемые?
  \end{enumerate}
\end{questionsbox}
